\section{Derivation of the mode-dependent Hamiltonian}\label{app:Hk_derivation}

In this Appendix, we derive the simplest case of the ideal one-bath scenario, obtaining \cref{eq:momentum-space-hamiltonian} as a result. It is worth noting that additional baths or environments can be incorporated by extending the resulting matrix, with each extra bath increasing the size by 2 rows and columns.

We begin our derivation at the operator level, performing a Fourier transform of the operators as follows:
\begin{align}
    \hat{H}_S &= \sum_{n=0}^{N-1} \cos\theta(\hat{a}_n^{\dag}-\hat{a}_n)(\hat{a}_{n+1}^{\dag}+\hat{a}_{n+1}) \nonumber\\
	          &\quad + \sin\theta(\hat{a}_n\hat{a}_n^\dag-\hat{a}_n^\dag \hat{a}_n), \label{eq:Hs_real_space}\\
    \hat{a}_n &= \frac{1}{\sqrt{N}}\sum_{k=-\frac{N}{2}}^{\frac{N}{2}-1}e^{-i\frac{2\pi k n}{N}}\hat{a}_k = \frac{1}{\sqrt{N}}\sum_{k=-\frac{N}{2}}^{\frac{N}{2}-1}e^{-i\phi_k n}\hat{a}_k, \label{eq:fourier_transform}
\end{align}
where $\hat{a}_n$ ($\hat{a}_n^\dag$) is the fermionic annihilation (creation) operator at site $n$, $N$ is the total number of sites, $\theta$ is a parameter of the system Hamiltonian, and $\phi_k = \frac{2\pi k}{N}$ is the phase factor associated with momentum mode $k$.

Substituting \cref{eq:fourier_transform} into \cref{eq:Hs_real_space} and simplifying yields:
\begin{align}
    \hat{H}_S &= \sum_k \cos\theta e^{-i\phi_k}(\hat{a}_k^\dag \hat{a}_{-k}^\dag+\hat{a}_k^\dag \hat{a}_{k})+\sin\theta\hat{a}_k\hat{a}_k^\dag + \text{h.c.} \nonumber\\
	          &= \sum_{k\geq 1}-2i\cos\theta\sin\phi_k[\hat{a}_k^\dag \hat{a}_{-k}^\dag+\text{h.c.}] \nonumber\\
	          &\quad +2(\cos\theta\cos\phi_k-\sin\theta)[\hat{a}_k^\dag \hat{a}_k+\hat{a}_{-k}^\dag \hat{a}_{-k}] \label{eq:Hs_momentum_space}
\end{align}